\documentclass[11pt,a4paper]{article}

\usepackage[spanish,es-tabla]{babel}
\usepackage[margin=2cm]{geometry}
\usepackage{amsmath}
\usepackage{booktabs}
\usepackage{graphicx}
\usepackage[table]{xcolor}
\usepackage{microtype}
\usepackage{enumitem}
\usepackage{titlesec}
\usepackage{fancyhdr}
\usepackage[hidelinks]{hyperref}
\setlist{nosep,leftmargin=*}
\setlength{\parindent}{0pt}
\setlength{\parskip}{5pt}
\definecolor{ink}{HTML}{13263A}
\definecolor{teal}{HTML}{0F766E}
\definecolor{coral}{HTML}{D95B43}
\definecolor{slate}{HTML}{52606D}
\definecolor{mist}{HTML}{EEF4F5}
\color{ink}
\titleformat{\section}{\Large\bfseries\color{ink}}{\thesection}{0.6em}{}
  [\vspace{2pt}\color{teal}\titlerule]
\titlespacing*{\section}{0pt}{12pt}{6pt}
\pagestyle{fancy}
\fancyhf{}
\lhead{\footnotesize\color{slate}Sistema inteligente adaptativo}
\rhead{\footnotesize\color{slate}UTEC · Proyecto 1}
\cfoot{\footnotesize\color{slate}\thepage}
\setlength{\headheight}{14pt}
\newcommand{\figdir}{figures/}

\title{\vspace{-1.2cm}\color{ink}\textbf{Sistema inteligente adaptativo para detección de fraude}\\
\large\color{slate}Evaluación temporal del dataset IEEE-CIS Fraud Detection}
\author{Edgar Chambilla\\Proyecto 1: Planificación y Toma de Decisiones en IA}
\date{Septiembre de 2026}

\begin{document}
\maketitle
\vspace{-0.5cm}

\begin{abstract}
Este trabajo desarrolla un sistema de clasificación para detectar fraude financiero en datos cuyo
comportamiento cambia con el tiempo. Parte de la línea base del informe técnico anterior, que integra
las tablas de transacciones e identidad, normaliza variables temporales y entrena XGBoost. La extensión
añade agregaciones causales por cliente implícito, ventanas expansivas y deslizantes, monitoreo de drift
con PSI y decisiones basadas en costos. La evaluación usa 590\,540 transacciones y bloques futuros de
30 días. La ventana expansiva alcanza un AUC promedio de 0,9163; la ventana de 30 días llega a 0,9039.
En esta corrida, conservar más historia funciona mejor que descartar datos de forma agresiva.
\end{abstract}

\section{Problema, objetivos y restricciones}

El problema es clasificar cada transacción como legítima o fraudulenta antes de autorizarla. El dataset
IEEE-CIS contiene 590\,540 transacciones, 20\,663 fraudes (3,50\,\%) y una tabla de identidad
disponible solo para una parte de los casos. La variable \texttt{TransactionDT} ordena aproximadamente seis
meses de actividad, por lo que una validación aleatoria produciría fuga de información.

El objetivo es maximizar la detección futura manteniendo controlados los falsos positivos. Las
restricciones son los datos tabulares ofuscados, el desbalance 1:27, la ausencia de un identificador
real de cliente, la memoria limitada de una RTX 2060 de 6 GB y el retraso en la disponibilidad de las
etiquetas. El protocolo no usa submisiones a Kaggle.

\textbf{Métricas.} Se reportan AUC y PR-AUC para ranking, precision, recall, F1 y balanced accuracy
para clasificación. Para decisiones se usa un umbral 0,20, costo 1 para falso positivo y costo 10
para falso negativo. También se reportan tiempo de entrenamiento y PSI.

\section{Diseño del sistema}

\begin{figure}[ht]
\centering
\includegraphics[width=\textwidth]{\figdir fig_arquitectura_adaptativa.pdf}
\caption{Arquitectura del sistema. El monitoreo alimenta la selección de ventana y el reentrenamiento.}
\end{figure}

El flujo recibe las tablas transaccional y de identidad, valida el orden temporal, calcula variables
causales, produce un score y lo convierte en una acción: aprobar, enviar a revisión o bloquear. Las
transacciones y decisiones se registran para medir desempeño y drift. La tabla transaccional describe
el evento de pago; la tabla de identidad aporta información sobre el dispositivo y la red.

\section{Preparación temporal y features}

Los valores faltantes se conservan como señal; las variables categóricas se codifican con un
ordinal encoder capaz de manejar categorías nuevas. Se derivan día relativo, hora, día de semana,
mes calendario y versiones normalizadas de $D4$, $D5$, $D10$, $D11$ y $D15$.

El cliente implícito se aproxima como:
\[
 UID = card1 \_ addr1 \_ \lfloor TransactionDT/86400-D1 \rfloor.
\]
El UID solo se utiliza como llave. Para cada fila se calculan frecuencia y medias históricas de monto,
variables temporales y conteos. Durante el entrenamiento se usa la suma acumulada anterior a la fila;
durante la evaluación se usa únicamente el histórico anterior al bloque. Así se elimina la fuga
transductiva del benchmark original.

\section{Modelado y protocolo experimental}

Se comparan regresión logística, Random Forest y XGBoost. La comparación auxiliar usa el mismo corte
temporal de la muestra de desarrollo. El experimento principal usa XGBoost, que fue el modelo con mejor
resultado en el informe base.

La evaluación adaptativa reserva los primeros 60 días para entrenamiento y evalúa cuatro bloques
futuros de 30 días. Se comparan:

\begin{itemize}
  \item histórico expansivo: conserva todos los días anteriores;
  \item ventana deslizante de 30, 60 o 90 días: olvida observaciones antiguas;
  \item PSI sobre monto, columnas $D$ y conteos $C$; PSI máximo mayor que 0,20 activa una alarma;
  \item reentrenamiento por bloque, con posibilidad de cambiar la ventana cuando aumente el costo.
\end{itemize}

\section{Resultados}

La tabla resume la corrida completa sobre las 590\,540 transacciones de entrenamiento, con XGBoost de
180 árboles y umbral 0,20.

\begin{table}[ht]
\centering
\small
\caption{Promedios de los cuatro bloques futuros.}
\begin{tabular}{lrrrrr}
\toprule
Ventana & AUC & PR-AUC & F1 & Recall & PSI máx.\\
\midrule
\rowcolor{mist}
Expansiva & 0,9163 & 0,5807 & 0,5575 & 0,5131 & 0,0403\\
30 días   & 0,9039 & 0,5630 & 0,5481 & 0,4786 & 0,0133\\
\rowcolor{mist}
60 días   & 0,9139 & 0,5759 & 0,5585 & 0,5045 & 0,0238\\
90 días   & 0,9156 & 0,5829 & 0,5635 & 0,5151 & 0,0295\\
\bottomrule
\end{tabular}
\end{table}

\begin{figure}[ht]
\centering
\includegraphics[width=\textwidth]{\figdir fig_adaptacion_metricas.pdf}
\caption{Evolución temporal de AUC, F1 y PSI. El segundo bloque muestra una degradación del modelo expansivo, aunque el PSI no supera el umbral alto.}
\end{figure}

La ventana expansiva alcanza el mejor AUC promedio. La ventana de 30 días presenta el peor resultado,
lo que indica que descartar historia puede eliminar patrones útiles. El rendimiento también cambia entre
bloques: el AUC expansivo baja de 0,9175 en el primero a 0,9071 en el segundo, mientras el costo de
decisión sube de 18\,059 a 21\,009. El sistema vigila el PSI y las métricas de desempeño cuando llegan
las etiquetas retrasadas.

\begin{table}[ht]
\centering
\small
\caption{Comparación auxiliar de modelos en el mismo corte temporal de desarrollo.}
\begin{tabular}{lrrr}
\toprule
Modelo & AUC & F1 & Recall\\
\midrule
\rowcolor{mist}
Regresión logística & 0,8022 & 0,0940 & 0,9686\\
Random Forest & 0,9000 & 0,1588 & 0,9186\\
\rowcolor{mist}
XGBoost & 0,9242 & 0,6163 & 0,5831\\
\bottomrule
\end{tabular}
\end{table}

\paragraph{El costo de la causalidad, medido.} El script \texttt{src/eval\_05\_integrated.py}
cuantificó la brecha transductiva sobre el sistema integrado del notebook 05. La variante protocolo
reproduce sus AUC (0,9340/0,9456 en holdout, 0,9347/0,9509 en OOF) y añade métricas de decisión en
umbrales 0,50, 0,19 y calibrado. La variante causal calcula las agregaciones del UID solo con el
pasado: entrenamiento en $[0\%,60\%)$, calibración en $[60\%,75\%)$ y evaluación en el futuro; el AUC
del holdout baja de 0,9456 a 0,9175. El umbral calibrado $\tau^{*}=0{,}108$ sostiene F1 0,60 en el
bloque futuro.

\begin{table}[ht]
\centering
\small
\caption{Sistema integrado 05 + magic UID, holdout causal 75/25 (\texttt{results/threshold\_metrics\_05\_causal.csv}).}
\begin{tabular}{lrrr}
\toprule
Umbral & Recall & Precisión & F1\\
\midrule
\rowcolor{mist}
0,50 & 0,3218 & 0,9308 & 0,4782\\
0,19 & 0,4200 & 0,8903 & 0,5707\\
\rowcolor{mist}
$\tau^{*}=0{,}108$ (calibrado) & \textbf{0,4602} & 0,8651 & \textbf{0,6008}\\
\bottomrule
\end{tabular}
\end{table}

\section{Despliegue, riesgos y adaptación}

En un despliegue por lotes diario o mensual, las nuevas transacciones pasan por validación, cálculo de
features históricas y predicción. El modelo se actualiza al cerrar cada bloque o cuando se detecta
drift. Los casos cercanos al umbral se envían a revisión humana; el sistema no debe bloquear de forma
irreversible sin auditoría. La incertidumbre se aproxima mediante la distancia al umbral y puede
mejorarse con calibración posterior.

Los principales riesgos son la fuga temporal en las agregaciones, los cambios en la cobertura de
identidad, las nuevas campañas de fraude, el costo desigual de los errores y la degradación silenciosa.
Para reducirlos se necesitan registros, versionado del modelo, rollback, revisión de alertas y
evaluación por período.

\section{Conclusiones}

\begin{sloppypar}
El informe base mostró que la ingeniería de entidades aporta más que cambiar de algoritmo. La extensión
adaptativa lleva esa evidencia a un flujo operativo: calcula features causales, evalúa bloques futuros y
compara estrategias de olvido. En esta ejecución, la ventana expansiva supera a la de 30 días. La
selección de ventana debe apoyarse en métricas de desempeño y costo, y el umbral conviene recalibrarlo
por bloque sin usar el futuro: en el sistema integrado, la variante causal del holdout costó 0,028 de
AUC y sostuvo F1 0,60 con $\tau^{*}=0{,}108$. El código reproducible se encuentra en
\texttt{src/adaptive\_fraud.py} y \texttt{src/eval\_05\_integrated.py}; los resultados, en
\texttt{results/adaptive\_results.csv}, \texttt{results/threshold\_metrics\_05\_integrated.csv} y
\texttt{results/threshold\_metrics\_05\_causal.csv}.
\end{sloppypar}

\end{document}
